\documentclass[10pt]{article}
\usepackage[margin=0.75in]{geometry}
\usepackage{xcolor}
\usepackage{tcolorbox}
\usepackage{hyperref}
\usepackage{enumitem}
\usepackage{booktabs}

% Define colors
\definecolor{osaorange}{RGB}{220,115,38}
\definecolor{aiblue}{RGB}{30,144,255}

% Configure hyperref
\hypersetup{
    colorlinks=true,
    linkcolor=blue,
    urlcolor=blue
}

% Compact spacing
\setlength{\parindent}{0pt}
\setlength{\parskip}{6pt}

\begin{document}

% Orange header
\begin{tcolorbox}[colback=osaorange,colframe=black,boxrule=2pt,arc=0pt,left=6pt,right=6pt,top=8pt,bottom=8pt]
\centering
\textcolor{white}{\textbf{\large WEEK 7 ASSIGNMENT}}\\
\textcolor{white}{\textbf{H524 Introduction to Biostatistics - Fall 2025}}
\end{tcolorbox}

\vspace{4pt}

\noindent\textbf{Total Points:} 10\\
\textbf{Due:} Sunday 11:59 PM

\vspace{8pt}

\section*{Part A: Multiple Choice Questions (5 points)}

Answer the following 5 questions on nonparametric methods. Each question is worth 1 point.

\subsection*{Question 1: When to Use Nonparametric Tests (1 point)}

Which of the following is the BEST reason to use a nonparametric test instead of a parametric test?

\begin{enumerate}[label=\Alph*., leftmargin=20pt]
\item Nonparametric tests are always more powerful than parametric tests
\item The data do not meet the assumptions required for parametric tests (e.g., normality)
\item Nonparametric tests require larger sample sizes
\item Parametric tests cannot be used with continuous data
\end{enumerate}

\subsection*{Question 2: Sign Test vs. t-test (1 point)}

A researcher wants to test whether a new medication reduces pain scores in patients. Pre- and post-treatment pain scores are measured for 15 patients. The pain score differences are highly skewed. Which test is most appropriate?

\begin{enumerate}[label=\Alph*., leftmargin=20pt]
\item Paired t-test
\item Sign test
\item Independent samples t-test
\item ANOVA
\end{enumerate}

\subsection*{Question 3: Mann-Whitney U Test (1 point)}

The Mann-Whitney U test is the nonparametric alternative to which parametric test?

\begin{enumerate}[label=\Alph*., leftmargin=20pt]
\item Paired t-test
\item One-way ANOVA
\item Independent samples t-test (two-sample t-test)
\item Chi-square test
\end{enumerate}

\subsection*{Question 4: Kruskal-Wallis Test (1 point)}

A study compares median recovery times across four different treatment groups. The recovery time data are highly right-skewed. Which test is most appropriate?

\begin{enumerate}[label=\Alph*., leftmargin=20pt]
\item One-way ANOVA
\item Kruskal-Wallis test
\item Chi-square test
\item Two-sample t-test
\end{enumerate}

\subsection*{Question 5: Ranking Data (1 point)}

In the Wilcoxon rank-sum test (Mann-Whitney U), how are tied values handled?

\begin{enumerate}[label=\Alph*., leftmargin=20pt]
\item Ties are broken randomly
\item Tied values receive the average of the ranks they would have occupied
\item Tied values are excluded from the analysis
\item Ties are not allowed in nonparametric tests
\end{enumerate}

\newpage

\section*{Part B: Group Project Check-In - Team Plan (5 points)}

\begin{tcolorbox}[colback=aiblue!10,colframe=aiblue,boxrule=1pt,arc=2pt,left=6pt,right=6pt,top=6pt,bottom=6pt]
\textbf{GROUP PROJECT BEGINS!}
\end{tcolorbox}

\vspace{4pt}

\textbf{Read First:} See \texttt{Group\_Project\_Overview.pdf} posted on Canvas for complete project details.

\subsection*{Your Assignment}

\textbf{Group Leader submits} a team plan including:

\textbf{1. Team Roster}
\begin{itemize}[leftmargin=20pt]
\item Names of all group members
\item Contact information (email, phone)
\end{itemize}

\textbf{2. Communication Plan}
\begin{itemize}[leftmargin=20pt]
\item How will you communicate? (email, text, Slack, Discord, etc.)
\item Preferred meeting format (in-person, Zoom, hybrid)
\end{itemize}

\textbf{3. Meeting Schedule}
\begin{itemize}[leftmargin=20pt]
\item At least 2-3 scheduled meeting times for Weeks 8-10
\item Include dates and times
\end{itemize}

\textbf{4. Division of Responsibilities}
\begin{itemize}[leftmargin=20pt]
\item Who will handle initial tasks:
\begin{itemize}[leftmargin=15pt]
\item Read data dictionary
\item Set up R and load dataset
\item Contact AI tools (Claude, ChatGPT)
\item Use Copilot in RStudio
\item Draft presentation outline
\item Write report sections
\end{itemize}
\end{itemize}

\textbf{5. Preliminary Timeline}
\begin{itemize}[leftmargin=20pt]
\item When will you complete data analysis?
\item When will you finish AI comparison?
\item When will you draft presentation?
\item When will you finalize report?
\end{itemize}

\subsection*{Submission Format}

\textbf{Submit as:} Text entry OR file upload (PDF/Word)

\textbf{Length:} 1-2 pages

\textbf{Grading:} Completion-based (5 points for submitting reasonable plan)

\subsection*{Notes}

\begin{itemize}[leftmargin=20pt]
\item \textbf{Only the group leader submits} - all team members get the same grade
\item Check Canvas to see your group number and dataset assignment
\item Download your dataset from: Canvas → Files → Group\_Project\_Data → \texttt{group\_X} folder
\item See data dictionary in Canvas Files for variable descriptions
\item This is just a PLAN - you can adjust as you go!
\end{itemize}

\end{document}
